% Generated draft: requires mathematical and source review.
\documentclass[11pt]{article}
\usepackage[margin=25mm]{geometry}
\usepackage{fontspec}
\setmainfont{Latin Modern Roman}
\usepackage{amsmath,amssymb,mathtools,unicode-math}
\setmathfont{Latin Modern Math}
\usepackage{enumitem,needspace,xurl,hyperref}
\setlength{\parindent}{0pt}
\setlength{\parskip}{5pt}
\setlength{\emergencystretch}{3em}
\newcommand{\N}{\mathbb{N}}
\newcommand{\Z}{\mathbb{Z}}
\newcommand{\Q}{\mathbb{Q}}
\newcommand{\R}{\mathbb{R}}
\newcommand{\C}{\mathbb{C}}
\newcommand{\F}{\mathbb{F}}
\newcommand{\A}{\mathbb{A}}
\newcommand{\PP}{\mathbb{P}}
\newcommand{\E}{\mathbb{E}}
\newcommand{\eps}{\varepsilon}
\newcommand{\dd}{\,\mathrm{d}}
\begin{document}
\section*{UnsolvedMath ID: 30004174}
\textbf{Source problem code:} catalogue entry
\subsection{Definitions and mathematical statement}
For the supplied permanent UnsolvedMath record \(30004174\), let \(X\) be a K-polystable Fano variety. The ground field and permitted singularities of \(X\) are not specified. For a positive integer \(l\) such that \(-lK_X\) is Cartier, write
\[
V_l=H^0\bigl(X,\mathcal O_X(-lK_X)\bigr).
\]
A nonzero section \(s_D\in V_l\) defines an effective divisor \(D=\operatorname{div}(s_D)\in|-lK_X|\); its scalar class is \([s_D]\in\mathbb P(V_l)\). This explains the usual section--divisor interpretation of the source notation \(H^0(X,K_X^{-l})\), but the source does not specify the allowed integers \(l\) or explicitly state the divisor hypothesis.

K-polystability and log K-polystability refer to the corresponding test-configuration stability notions: the relevant Donaldson--Futaki invariant is nonnegative, with equality permitted only for product configurations, subject to the conventions and admissibility hypotheses of the chosen setting. GIT stability means geometric invariant theory stability for a specified group action and linearization. The source names the natural \(\operatorname{Aut}(X)\)-action on plurianticanonical sections, but does not specify the linearization or whether stability is taken in \(V_l\) or \(\mathbb P(V_l)\).

The supplied question asks whether, for every such K-polystable Fano variety \(X\) and for sufficiently large \(\beta\), the equivalence
\[
(X,(1-\beta)D)\text{ is log K-polystable}
\quad\Longleftrightarrow\quad
\text{a section defining }D\text{ is GIT-stable under }\operatorname{Aut}(X)
\]
holds. This is a statement of the intended target, with unresolved parameters, rather than a fully quantified conjecture: the source gives neither the range of \(\beta\), the quantification over \(l\) and \(D\), nor the permitted dependence of the threshold for \(\beta\). In particular, interpreting ``sufficiently large'' as approaching \(1\) from below would require an additional convention absent from the record.

There is also an unresolved normalization issue. If \(D\in|-lK_X|\) is used without rescaling, then
\[
-\bigl(K_X+(1-\beta)D\bigr)
\sim_{\mathbb R}\bigl(1-l(1-\beta)\bigr)(-K_X).
\]
Thus the usual effective-boundary log Fano interpretation requires appropriate coefficient, singularity, and positivity hypotheses. Replacing \(D\) by \(D/l\) would change the displayed pair and is not authorized by the supplied statement. Finally, the source explicitly compares log K-\emph{polystability} with GIT \emph{stability}; it does not resolve whether this distinction is intentional.
\subsection{Short English statement}
The question asks whether adding a sufficiently small divisor boundary to a K-polystable Fano variety makes log K-polystability equivalent to GIT stability of the section defining that divisor. The phrase ``sufficiently small boundary'' would correspond to \(\beta\) near \(1\) from below, an interpretation that the source does not explicitly specify. Divisor normalization, parameter quantifiers, and the precise GIT notion require clarification before a fully precise conjecture can be stated.
\subsection{Sources}
\begin{itemize}
\item Supplied local UnsolvedMath registry record, permanent ID \(30004174\), titled ``Log K-Stability versus Geometric Invariant Theory,'' source problem code \texttt{catalogue entry}, in the Oberwolfach Reports Open Problems collection. No individual proposer is supplied; the proposed year is recorded as \(2019\).
\item Supplied report citation: \emph{Differentialgeometrie im Großen (2020)}, Oberwolfach Reports, DOI \texttt{10.4171/owr/2019/30}, pp.\ 1791--1839. Original OWR report: \url{https://doi.org/10.4171/owr/2019/30}. The supplied background describes these reports as records of mathematical workshops and open-problem discussions.
\item Zhou, \emph{Log K-stability of GIT-stable divisors}: \url{https://arxiv.org/abs/2102.07458}.
\item Literature assessment dated exactly \texttt{2026-08-22}: supplied dated curation, not independently verified. The record labels the problem ``partially solved'' and states that Zhou proves the GIT--log K-stability equivalence for K-polystable Fano varieties under additional log-canonicity and small-coefficient hypotheses, while the unrestricted formulation for arbitrary plurianticanonical divisors remains unresolved. The background identifies this assessment as dated literature triage derived from the supplied OWR curation and requests independent verification of time-sensitive status and cited claims.
\item The registry supplies the statement label \texttt{corrected\_verified} and reports a manual semantic, typographic, and LaTeX review. These are supplied editorial assertions; neither the sources nor their mathematical claims have been independently checked in preparing this draft.
\end{itemize}
\end{document}
